% 1995.03.20  LaTeX_edit(english version) by Computer Institute of Japan.
\documentstyle[epsf,art10]{article}

\textheight=25cm
\textwidth=16cm
\topmargin=-1.0cm
\oddsidemargin=0cm


\begin{document}
\baselineskip=19pt

\title{GUI operation manual}
\author{}
\date{}

\pagestyle{empty}
\pagestyle{plain}

\maketitle

%\tableofcontents

%\newpage

\makeatletter
\def\section{\@startsection {section}{1}{\z@}{3.5ex plus 1ex minus 
.2ex}{2.3ex plus .2ex}{\Large\bf}}
\def\subsection{\@startsection{subsection}{2}{\z@}{3.25ex plus 1ex minus 
.2ex}{1.5ex plus .2ex}{\large\bf}}
\def\subsubsection{\@startsection{subsubsection}{3}{\z@}{3.25ex plus 
1ex minus .2ex}{1.5ex plus .2ex}{\normalsize\bf}}
\makeatother

\normalsize

\section{Overview}

This manual explains how to debate using the GUI (Graphical User Interface) of "HELIC-II".
Conceptual explanation of a debate function is given in '\ref{Debate model}'. 
It is recommended to read this chapter before using GUI.
Terminology used in this manual are explained where they first appear.
However, it is strongly recommended to read through the language manual for further and systematic understanding.

\section{Basic knowledge}

\subsection{Debate model}\label{Debate model}

A debate function of "HELIC-II" materializes a debate model between two agents called a prosecutor and defendant.

The object of the prosecutor agent is to substantiate an argument to justify a self goal, and the object of the pleading agent is to substantiate an effective counter argument to the prosecutor's argument.
A premise of this debate model is as follows:

\begin{itemize}

\item Each agent has a different knowledge (viewpoint, rule set, etc.).
\item Only the prosecutor agent has a goal to achieve in a debate.
This goal, represented as an H term\footnote{\ For details, refer to the language manual.}, may include a variable.
A defendant rebuts passively to the argument sent from the prosecutor.
\item Claims from both sides must use arguments in a form of a proof tree based on facts.
\item The comparison between arguments is determined by the priority
  relation between the TopRule of arguments\footnote{\ Comparing without evaluating the Sub-Argument causes a different result from the Defeasible Reasoning.}.

\end{itemize}

A debate is proceeded by claiming an argument to the opponent and noticing that the argument is advantageous over the other.

% +----+  	---->  		+----+
% |$B8!;!(B|     $BO@Ah%W%m%H%3%k(B 	|$BHo9p(B|
% +----+  	<----  		+----+
%   |		  	  	  |
%   V		  	  	  V
% +===========+	       		+===========+
% |$B2ACM4Q(B1    |	       		|$B2ACM4Q(B2    |
% |$B%k!<%k=89g(B1|	       		|$B%k!<%k=89g(B2|
% +===========+	       		+===========+
% 
% 		[ $BO@Ah%b%G%k(B ]

Each agent sends a message called a debate protocol to the opponent to present the argument.
The following are the different kinds of debate protocols:

\begin{enumerate}

\item Claim argument (pose)\\
         Claims an argument.  It is used in a case of doing a claim or rebuttal.
\item Priority notice (notice)\\
         Notices that the self argument is dominant.
\item Cancel argument (cancel)\\
         Cancels a claimed argument.
\item End debate (end)\\
         Notices that the debate on the opponent's argument has been finished.

\end{enumerate}

\section{Structure of HELIC-II GUI}

\subsection{Constitution}

HELIC-II GUI is composed of two Unix Processes.
One is the inference engine module written by KL1 and the other is the GUI portion written in C.
These two Unix Processes are tied together by a socket communication.
To boot "HELIC-II" the inference engine must be first started, followed by the  GUI.
For instruction on how to boot, refer to '\ref{Activation method}'.

\subsection{State transition during a GUI execution}

The GUI which runs the debate function has the following three states, and the commands that may be used depend on these three states.
It is in the initial state at the time of booting.

%	   	  debate,rbr		      start
%+--------+  	--------> 	+--------+  --------> 	+--------+
%|$B=i4|>uBV(B|		  	|$B@_Dj>uBV(B|              |$BO@Ah>uBV(B|
%+--------+	<-------- 	+--------+  <-------- 	+--------+
%		  quit			      end

\begin{enumerate}
\item Initial state:

\begin{enumerate}
\item Use possibility commands
\begin{itemize}
\item Knowledge base load relation
\item Debate agent role setting
\end{itemize}
\item State transition command\\
{\sl Debate} or {\sl RDB} changes the setting state.
\end{enumerate}

\item Setting state: 

\begin{enumerate}
\item Use possibility commands
\begin{itemize}
\item Viewpoint binding
\item Goal binding
\item Use unit binding
\item Argument browser command\footnote{\ Available only when an
    argument has already been displayed, and the debate has been finished.}
\end{itemize}
 
\item State transition commands\\
{\sl Quit} returns it to the initial state, and {\sl Start} begins the debate.
\end{enumerate}

\item Debate state:

\begin{enumerate}
\item Use possibility commands
\begin{itemize}
\item Debate relative command
\item Argument browser command
\end{itemize}
\item Condition transition command\\
{\sl End} returns it to the setting state.
\end{enumerate}

\end{enumerate}

The following processing is required to transit from the initial state to the setting state:
\begin{itemize}
\item Loading a knowledge file.
For details, refer to '\ref{To load a knowledge}'
\item Setting debate agents.
For details, refer to '\ref{To set roles of a debate}'.
\end{itemize}

The following processing is required to go from a setting state to a debate state:
\begin{itemize}
\item Set a use knowledge.
For details, refer to '\ref{To set a use knowledge}'.
\item Set a viewpoint.
For details, refer to '\ref{To set a viewpoint}'.
\end{itemize}

\section{Activation method}\label{Activation method}

To start "HELIC-II" GUI, the inference engine and the GUI must be started from the Unix shell.
As mentioned earlier, the inference engine must be started before the GUI. 

\subsection{Start commands}

Commands in this section, {\sl guidbt} and {\sl gui}, are available in a "bin" directory under the installation directory.
Refer to the installation manual (INSTALL) on how to create these commands.

\subsubsection{Inference engine}

{\bf Command:}
\begin{quote}
\verb+guidbt [-port PORTNO]+
\end{quote}

{\bf Argument:}
\begin{quote}
\begin{tabbing}
PORTNO \ \ \= \kill
\verb+PORTNO+ \> The Port number of socket.\\
\end{tabbing}
\end{quote}

{\bf Explanation:}
\begin{quote}
To start the inference engine, specify the port number of the socket that will connect the inference engine with GUI Process by an argument or an environment variable HELIC\_HOST\_PORT.
The port number must be the same as the port number specified for the GUI.

The following shows how to specify a port number which is 9900 in this example:

\begin{verbatim}
% setenv HELIC_HOST_PORT 9900
\end{verbatim}

\end{quote}

{\bf Remarks:}
\begin{quote}
The port number will be automatically set to 9000 when it is not specified with an environment variable.
This command ends automatically by ending the GUI.

\end{quote}

\subsubsection{GUI}

{\bf Command:}
\begin{quote}
\verb+gui [-host HOSTNAME] [-port PORTNO]+
\end{quote}

{\bf Argument:}
\begin{quote}
\begin{tabbing}
HOSTNAME \ \ \= \kill
\verb+HOSTNAME+ \> The host name of the machine that started the inference engine Process.\\
\verb+PORTNO+ \> The port number of the socket.\\
\end{tabbing}
\end{quote}

{\bf Explanation:}
\begin{quote}
Start GUI Process.
The inference engine must have already been started at this point.
Specify the port number of the socket and the host name of the machine the inference engine is currently running on.
This specification may also be made by the environment variable.

The host name of the machine which the inference engine Process is currently running is specified by the environmental variable HELIC\_HOST\_NAME.
The port number of the socket to be connected to the inference engine Process is specified by the environmental variable HELIC\_HOST\_PORT.
This port number must be the same as the port number specified when the inference engine first started. 

The specification of an argument has priority when both methods are used to specify a port number and/or a host name.  When the two are not specified, an error window appears and the GUI ends.

The initial window of the GUI is then presented on the screen.

{\bf Related files:}

\begin{quote}
~\verb+/.Xdefaults+
\end{quote}
\end{quote}

\newpage

\section{Operation in the initial window}

This window first opens when the "HELIC-II" is started.
A knowledge base must be loaded and debate agents must be assigned in this window before opening the debate window.
For details on how to start this window, refer to '\ref{Activation method} Activation method'.
In a standard environment, simply enter the following:
\begin{verbatim}
	% setenv HELIC_HOST_PORT 9900
	% setenv HELIC_HOST_PORT here
	% debate &
	% gui
\end{verbatim}

\begin{figure}[htb]
\begin{center}
\epsfile{file=./initwin.ps,height=6cm}
\end{center}
\caption{Initial window}\label{Initial window}
\end{figure}

\begin{itemize}
\item Browsing area
\begin{description}
\item[Loading Status:] Suggests the knowledge base file that must be loaded in order to start the debate.
\item[Role Assigment Status:] Displays which debate agent the user or the system is taking charge.
\end{description}
\item Action area
\begin{description}
\item[Load button:] Loads a knowledge base file.
Opens the file load window (refer to '\ref{To load a knowledge}').
\item[Role button:] Determines roles of debate agents.
\item[Debate button:] Opens a debate window.
This button does not work until the load and the role setting end.
\item[Quit button:] Closes all windows and ends "HELIC-II".
\end{description}
\end{itemize}

\subsection{To end HELIC-II}
Select {\sl Quit} in the initial window.

\subsection{To load a knowledge}\label{To load a knowledge}
Select {\sl Load} in the initial window and open the file load window.

\begin{figure}[htb]
\begin{center}
\epsfile{file=./loadwin.ps,height=10cm}
\end{center}
\caption{File load window}\label{File load window}
\end{figure}

\begin{itemize}

\item File type selecting area
\begin{description}
\item[Type:] Select a type file.
Displays *.t or *.type file on the List field.
\item[Rule:] Select a rule file.
Displays *.r or *.rule file on the List field.
It may be selected only when a type file is loaded.
\item[View:] Select a viewpoint file.
Displays *.v or *.view file on the List field.
It may be selected only when a rule file is loaded.
\item[Query:] Select a goal file.
Displays *.g or *.goal file on the List field.
It may be selected only when a viewpoint file is loaded.
\item[Command:] Select a command file.
Displays *.c file on the List field.
\end{description}

\item File selecting area
\begin{description}

\item[Filter:] Regular expression of the file to be displayed on the Filter File. 
\item[Directory:] Displays a list of directories under the current directory.
\item[List:] Displays a list of knowledge base files of the type selected in the file type selecting area in the current directory.
\item[Loading File:] Displays file names that may be loaded.
It may be input directly from a keyboard.

\end{description}

\item Action area
\begin{description}
\item[Load button:] Orders an inference engine module to load the selected knowledge base file.
\item[List button:] Opens a list display window of knowledge base files loaded.

\begin{figure}[htb]
\begin{center}
\epsfile{file=./listwin.ps,height=3cm}
\end{center}
\caption{Load file list window}\label{Load file list window}
\end{figure}

\item[Done button:] Ends a load operation.Closes the file load window.

\end{description}

\end{itemize}

\subsubsection{To move a directory}
Double click the directory displayed in the Directory field of the file load window with the mouse.

\subsubsection{To load a type file}
Click {\sl Type} at the upper part of the file load window to display *.t or *.type on the List field.
Double click the file to load. 

\subsubsection{To load a rule file}

Click {\sl Rule} at the upper part of the file load window to display *.r or *.rule on the List field.
Double click the file to load. 

\subsubsection{To load a viewpoint file}

Click {\sl View} at the upper part of the file load window to display *.v or *.view on the List field.
Double click the file to load. 

\subsubsection{To load a goal file}

Click {\sl Goal} at the upper part of the file load window to display *.g or *.goal on the List display.  
Double click the file to load.

\subsubsection{To run a command file}

Click {\sl Command} at the upper part of the file load window to display *.c file on the List field.  
Double click the file to execute.

\subsubsection{To see a list of files loaded}

Select {\sl List} in the file load window and then the load file list window will be opened.

\subsection{To set roles of a debate}\label{To set roles of a debate}

To set roles of the debate, select {\sl Role} in the initial window and then the role setting window will be opened.

\begin{figure}[htb]
\begin{center}
\epsfile{file=./rolewin.ps,height=5cm}
\end{center}
\caption{Role setting window}\label{Role setting window}
\end{figure}

\begin{itemize}
\item Role setting area \\
Set the role of a prosecutor at the left and the role of a defendant at the right.
\begin{description}
\item[System:] The system will take charge of the given debate agent (prosecutor or defendant).  The debate will be proceeded automatically.
\item[User:] The user will take charge of the given debate agent.  The user is responsible for the operation needed for the debate to progress.
\end{description}

\item Action area
\begin{description}
\item[Done button:] Registers roles to a system.
Closes the role setting window.

\end{description}

\end{itemize}

\subsubsection{To set a user take charge of a prosecutor (defendant)}
Click {\sl User} at the left (right, in the case of a defendant) side of the role setting window.

\subsubsection{To set a prosecutor (defendant) to be done automatically}
Click {\sl System} at the left (right, in the case of a defendant) side of the role setting window.

\subsubsection{To register a setting}
Select {\sl Done} in the role setting window.

%%%%%%%%%%%%%%%%%%%%%%%%%%%%%%%
%
%	$BO@Ah%&%$%s%I%&(B
%
%%%%%%%%%%%%%%%%%%%%%%%%%%%%%%%

\section{Operation in the debate window}
\begin{quotation}
Two debate windows are opened simultaneously to represent the prosecutor and defendant.

Various setting and operations done in the prosecutor window will apply toward the prosecutor.

To open a debate window, a type, viewpoint, goal, and rule must all be loaded, and the settings of the prosecutor and defendant completed. A pop up window from a debate window is an exclusive control window; while being operated, it will prevent the debate window from receiving any control.
The debate window becomes sensitive only when it is its turn to speak.

\end{quotation}

\begin{figure}[htb]
\begin{center}
\epsfile{file=./debatewin.ps,height=12cm}
\end{center}
\caption{Debate window}\label{Debate window}
\end{figure}

\begin{itemize}

\item Debate setting part\\
The following buttons becomes insensitive during a debate:
\begin{description}
\item[View button:] Sets a viewpoint. 
\item[Unit button:] Sets a use knowledge. 
\item[Goal button:] Sets a goal. 
\item[Start button:] Starts a debate. 
\item[Quit button:] Closes a debate window. 
\end{description}

\item Debate setting browser part
\begin{description}
\item[Goal area:] Displays the goal that is set. 
\item[View area:] Displays the viewpoint that is set. 
\end{description}

\item The following buttons are insensitive except during a debate in
  the debate operation part, and only the executable commands become sensitive at such time.
\begin{description}
\item[Pose button:] Claims the argument displayed in the logic display area to the opponent. 
\item[Notice button:] Gives a priority notification to the opponent showing that its current argument displayed in the show argument area has higher priority than the opponent's argument.
\item[Cancel button:] Withdraws the claim displayed in the show argument area.
\item[Finish button:] Informs that the counter argument to the opponent's argument displayed in the show argument area has ended.
\item[Justify button:] Justifies an argument.
\item[Depth button:] Sets the number of steps to plans ahead for an opponent's counter argument before generating a claim or counter argument.
\item[End button:] Forces a debate to end.
\end{description}

\item Debate state browser part
\begin{description}
\item[Debate history area:] Displays a list that shows the states of the present debate.
\item[State show area:] Display the state of a debate.
The following are examples of states displayed:
"The opponent has claimed the following counter argument".
"The opponent has noticed that the following claim has priority under the human right esteem".
\item[Issue point display area:] Displays the issue point of the opponent's claim which the argument currently displayed is counter arguing.
\item[Show argument area:] Displays arguments.
Displays a Unified conclusion (intermediate proposition) of rules with an indentation, depending on the reference relation.
Each rule composes the proof tree of an argument.
\end{description}

\item{Argument browser operation part}
\begin{description}
\item[Backward button:] Displays the argument which the counter argument currently displayed in the show argument area is trying to override.
\item[Forward button:] Returns to the argument displayed immediately
  before selecting {\sl Backward}, in the case where {\sl Backward} was used.
Returns to one of the counter arguments to the argument displayed in the show argument area.
\item[Next button:] Advances to the next counter argument, if it exists, disputing on the same issue point number as the argument currently displayed in the show argument area.
\item[Previous button:] Advances to the previous counter argument, if it exists, disputing on the same issue point number as the argument currently displayed in the show argument area.
Displays in the opposite order as the {\sl Next} button.
\item[Focus button:] Jumps to the argument (Focus argument) currently focused in the debate.
\item[ClaimList button:] Retrieves its or opponent's past claims.
\item[Jump button:] Displays the specified argument when a claim identification is input.
\item[Explain button:] Explains the argument currently displayed in a natural language.
\item[Form button:] Specifies the display form of the intermediate proposition displayed in the show argument area.
There are two display formats available; one is an expression by an H term variable and the other is an explanation in a natural language.
\end{description}

\end{itemize}

\subsection{To open a debate window}
Select {\sl Debate} in the initial window.
However, the role setting and loading of required files must be completed for the {\sl Debate} button to become sensitive.

\subsection{To close a debate window}
Select {\sl Quit} button.

\subsection{To set a viewpoint}\label{To set a viewpoint}
Select {\sl View} in the debate window and then a viewpoint setting window will be opened.

\begin{figure}[htb]
\begin{center}
\epsfile{file=./viewset.ps,height=6cm}
\end{center}
\caption{Viewpoint binding window}\label{Viewpoint binding window}
\end{figure}

\begin{description}

\item[View Candidate button:] Calls a viewpoint candidate selecting window and sets a viewpoint by selecting from a list of viewpoint candidates.
\item[Standpoint button:] Calls a standpoint selecting window and adds the selected standpoint to a viewpoint.
\item[Viewpoint display area:] Displays the content of the current viewpoint.
Displays the name of a standpoint in the order of higher priority.
\item[Order button:] Changes the order of a standpoint.
Opens a viewpoint order modification window.
\item[Done button:] Sets the specified viewpoint to a debate agent.
\item[Cancel button:] Clears a viewpoint setting work.
\end{description}

\subsubsection{To set a viewpoint by specifying a viewpoint candidate}
To specify a viewpoint candidate that is already available, select {\sl View Candidate} to open a view point candidate selecting window.
Double click a viewpoint candidate from the list displayed in the viewpoint candidate selecting window.
A list of standpoints of the selected viewpoint candidate is displayed in the viewpoint display area.
Select {\sl Done} to set the viewpoint to the debate agent, or use {\sl Standpoint} button or {\sl Order} button to edit.

\subsubsection{To add a standpoint to a viewpoint}
Select {\sl Standpoint} to open a standpoint selecting window, where every defined standpoint is listed.
Double click the name of the standpoint to be added.
The selected standpoint will be added at the bottom of the list of the viewpoint.

\subsubsection{To delete a standpoint from a viewpoint}
Double click the standpoint to delete in the viewpoint display area.

\subsubsection{To modify the order of a viewpoint}
Select {\sl Order} to open a viewpoint order modification window.
The current order is displayed to the left.
Double click standpoints in the order to be placed.
The standpoints will be displayed to the right in the new order.
After selecting all of the standpoints, the viewpoint order modification window closes automatically and the viewpoint will be displayed in the new order in the viewpoint display area.

\subsubsection{To clear a viewpoint}
Select {\sl Clear}.

\subsubsection{To register the modified viewpoint}
Select {\sl Done}.  The viewpoint displayed in the viewpoint display
area will be registered as the viewpoint of the debate agent.


\subsubsection{To quit viewpoint modification work}
Select {\sl Cancel}.
The modification work is cleared and the viewpoint modification window closes. 
No changes are made to the viewpoint.

\subsection{To set a use knowledge}\label{To set a use knowledge}

Select {\sl Unit} in the debate window to open a knowledge setting window.

\begin{figure}[htb]
\begin{center}
\epsfile{file=./unitset.ps,height=6cm}
\end{center}
\caption{Knowledge setting window}\label{Knowledge setting window}
\end{figure}

\begin{description}
\item[Unit Name:] Area to specify a meta-unit name for a use knowledge.
\item[Assigned Unit List:] List of meta-units specified.
\item[Done:] Registers a unit name input.
\item[Cancel:] Clears the knowledge setting work and closes the knowledge setting window.
\end{description}

\subsubsection{To specify a use knowledge}
Set the cursor to the field next to the unit name and input the unit name.
If there are no errors in the unit name, it will be added to the Assigned Unit List.
An error window is output when errors occur.

\subsubsection{To delete a use knowledge}
Double click the unit name to delete in the Assigned Unit List.
\subsubsection{To clear a use knowledge}
Select {\sl Clear}.

\subsubsection{To register a modified use knowledge}
Select {\sl Done}.
The list of unit names displayed in the Assigned Unit List will be registered as the use knowledge.

\subsubsection{To quit an use knowledge modification work} 
Select {\sl Cancel}.
The modification work is cleared and the use knowledge modification window closes. 
No changes are made to the viewpoint.

\subsection{To set a goal} 
Select {\sl Goal} in the debate window to open a goal binding window.

\begin{figure}[htb]
\begin{center}
\epsfile{file=./goalwin.ps,width=15cm}
\end{center}
\caption{Goal binding window}\label{Goal binding window}
\end{figure}

\begin{description}
\item[Goal List:] Displays goal candidates in H term expressions.
The goal candidate currently set is highlighted when the goal set window is opened. 
\item[Cancel button:] Clears the goal binding work, and closes the goal binding window.
\end{description}

\subsubsection{To specify a goal}
\medskip
Double click on the goal candidate displayed in the Goal List.  The goal set window closes automatically.

\subsubsection{To quit a goal setting work}
\medskip
Select {\sl Cancel}.
The goal setting window closes. 
No changes are made to the goal.

\subsection{To do a debate} 

\subsubsection{To start a debate}
Select {\sl Start} in the debate window.
The viewpoint and use knowledge of both the prosecutor and defendant, and the goal of the prosecutor must be set.
If the user is the prosecutor, a claim selection operation must be
made after {\sl Start}.
For details regarding this operation, refer to '\ref{To select a claim} To select a claim'.

\subsubsection{To finish a debate} 
Select {\sl End} in the debate window.

\subsubsection{To select a claim}\label{To select a claim} 
If the user is the prosecutor, a claim selecting window opens.
Select an initial claim, which will guide the prosecutor's claim, when the debate first starts, and when the opponent notices that the counter argument to the first argument has finished.

\begin{figure}[htb]
\begin{center}
\epsfile{file=./argslctwin.ps,height=6cm}
\end{center}
\caption{Claim selecting window}\label{Claim selecting window}
\end{figure}

A list of conclusions of the initial claim is displayed in this window.
Refer to '\ref{Argument list display of a viewpoint}Argument list display of a viewpoint' for the meaning of \verb+>0:!0:<1+.
Double click a conclusion of a claim from the list.
The claim selecting window closes and the selected claim (argument) will be displayed in the show argument area.

\subsubsection{To assert a claim to an opponent}
To assert an initial claim or a counter argument to the opponent, an argument needs to be sent to the opponent.
To do so, the following two must be done beforehand:
\begin{itemize}
\item Display the argument to send in the show argument area.\\
For details, refer to '\ref{Operating method of an argument browser} Operating method of an argument browser'.
\item The counter argument to be sent must be effective, that is, in the case of a prosecutor, the argument must have a dominant position, and in the case of a defendant, the argument must have a dominant position or be unrelated.\\
For details on dominant positions refer to '\ref{To make an argument have a dominant position} To make an argument have a dominant position'. 
\end{itemize}
If the above condition is met, the {\sl Pose} button of a debate operation part becomes sensitive.
Select {\sl Pose} to send and assert claims to the opponent and give the turn to speak to the opponent.

\subsubsection{To choose a counter argument}

A claim of the opponent will be displayed on the show argument area.
Select the issue point to do a counter argument.
The number to the left of the issue point indicates which issue point
to counter argue\footnote{\ For details, refer to '\ref{Argument display of a viewpoint} Argument display of a viewpoint'.}.
Double click the issue point to open a counter argument selecting window.

\begin{figure}[htb]
\begin{center}
\epsfile{file=./cargwin.ps,height=6cm}
\end{center}
\caption{Counter argument selecting window}\label{Counter argument selecting window}
\end{figure}

A list of conclusions of counter arguments which guides the opposite argument to the issue point is displayed in this window.
Refer to section '\ref{Argument list display of a viewpoint} Argument display of a viewpoint' for the meaning of \verb+>0:!0:<1+.
Double click a conclusion of a counter argument from the list.
The counter argument selecting window closes and the selected counter argument will be displayed in the show argument area.

\subsubsection{To give priority notice}\label{To give priority notice}

The opponent's counter argument can be ended by declaring a dominant (or unrelated in the case of a defendant) position of the original argument.
The following conditions must be met to do so:

\begin{itemize}
\item The argument to claim a dominant position is displayed in the show argument area.\\
Use the {\sl Backward} button, since the counter argument of the opponent is displayed.
\item The argument is not weaker than the opponent's counter argument.\\
For details on how to make an argument have dominant position, refer to '\ref{To make an argument have a dominant position} To make an argument have a dominant position'.
\end{itemize}
When these conditions are met, the {\sl Notice} button of a debate operation part becomes sensitive.
Select {\sl Notice} to open a standpoint selecting window.

\begin{figure}[htb]
\begin{center}
\epsfile{file=./spslctwin.ps,height=6cm}
\end{center}
\caption{Standpoint selecting window}\label{Standpoint selecting window}
\end{figure}

A list of standpoints or a combination of standpoints which has become the basis to judge the dominant position is displayed on the window.
Double click the standpoint to inform the opponent.
This completes the process of a priority notice.

However, when the defendant sends an unrelated argument, the standpoint selecting window does not open, and the turn to speak will change to the opponent.

\subsubsection{To inform priority over a priority notification from an opponent}

Display the counter argument that was claimed invalid in the debate display area, and send a priority notice to the opponent.
For details on how to send a priority notice, refer to section '\ref{To give priority notice} To give priority notice'.
The turn to speak will not be changed to the opponent.  Search for a different counter argument to the opponent's claim.

\subsubsection{To abandon a claim}

When an effective counter argument cannot be found, the claim must be abandoned.
Select {\sl Cancel} to abandon the claim.

\subsubsection{To end a counter argument}

Display the argument in the debate display area.\\
Select {\sl Finish} to end its counter argument.  This will demand another argument from the opponent.

\subsubsection{To make an argument have a dominant position}\label{To make an argument have a dominant position}

The current viewpoint can be reinforced to make the argument displayed have a dominant position.
Select {\sl Justify} to search for the standpoint to add, in order to make the argument displayed stronger than that of the opponent.
When such combination is found, the viewpoint change suggestion window opens.
Select a combination from the list displayed in this window.
The viewpoint is modified and the content of the show argument display is updated.

\begin{quote}
\begin{figure}[htb]
\begin{center}
\epsfile{file=./viewchgwin.ps,height=6cm}
\end{center}
\caption{Viewpoint change suggestion window}\label{Viewpoint change suggestion window}
\end{figure}

The list on the left shows the combination of standpoints that may be added.  Select by clicking on a standpoint .
The content of the viewpoint with the added standpoint is displayed in a list at the right.
Select {\sl Done} to register the modification of the viewpoint and to close the viewpoint change suggestion window.
\end{quote}

The change in viewpoint is done by retrieving predefined standpoints.  Therefore, a viewpoint to change the dominant position may not found.
Also, new viewpoints which contradict the priority notification already given will not become candidates.
Therefore, selecting {\sl Justify} does not necessarily retrieve standpoints to add.  In such case, an error window opens.
In such cases, find another argument or abandon the previous argument to proceed.

\subsubsection{To modify steps to plan ahead}

Claims and counter arguments are being generated with a frame of Defeasible Reasoning.  In Defeasible Reasoning, a counter argument to an argument is requested recursively.  Therefore, in order to decrease the amount of computing, the number of steps to plan ahead to an argument can be limited.\\
Select {\sl Depth} to open the plan ahead (foresight?) setting window.
Enter the number of steps to plan ahead.
Select {\sl Done} to complete.
Entering 1 indicates that one counter argument from the opponent can displayed.

\subsection{Operating method of an argument browser}\label{Operating method of an argument browser}

An argument browser is a function to display an argument in a show argument area using buttons of a debate operation part of a debate window.

\subsubsection{Argument display of a viewpoint}\label{Argument display of a viewpoint}

The following is an example of a display form of a general argument in a show argument area:
\begin{figure}[htb]
%\begin{center}
\begin{verbatim}
	general_penal_rule00
	weak
	             provable_guilty(a_object=take_away(agent=person[nam
	              take_away(agent=person[name=>"mary"],object=handba
	              related_crime(a_object=take_away(agent=person[name
	              punishable(a_object=take_away(agent=person[name=>"
	               satisfied_condition_of_crime(a_object=take_away(a
	                take_away(agent=person[name=>"mary"],object=hand
	                is_stealing(a_object=take_away(agent=person[name
	                 has_occupied(a_object=person[name=>"jane"],obje
	                 take_away(agent=person[name=>"mary"],object=han
	                 recognize(a_object=person[name=>"mary"],object=
	                  recognize(a_object=person[name=>"mary"],object
	                   mistake(agent=person[name=>"mary"],object=fai
	                  take_away(agent=person[name=>"mary"],object=ha
	   !  1(0)      with_criminal_intent(a_object=person[name=>"mary
	                 take_away(agent=person[name=>"mary"],object=han
	               not -punishable(a_object=take_away(agent=person[n
\end{verbatim}
%\end{center}
\caption{Example of an argument display}\label{Example of an argument display}
\end{figure}

The first line displays the name of the TopRule of the displayed argument.
The second line shows the priority relation between the TopRule of the
displayed argument and the rule which guides the intermediate
proposition of the parent argument which is refuted by the displayed argument.
{\bf Weak} means that the TopRule of the displayed argument has lower priority.
{\bf Strong} means that the TopRule of the displayed argument has higher priority.  {\bf Non-relation} indicates that there are no priority relations between the two.
Lines with a number added to the beginning mean an argument which is expressed in a form of a proof tree.
Each line is an intermediate proposition that was guided by an intermediate rule used in an argument.
All variables are expressed in the Unified form.
The following is another example to explain in details:
\begin{figure}[htb]
%\begin{center}
\begin{verbatim}

 Child information  intermediate proposition 
=====================================================================
       !  1(0)      with_criminal_intent(a_object=person[name=>"mary

\end{verbatim}
%\end{center}
\caption{Example of an intermediate proposition of an argument display}\label{Example of an intermediate proposition of an argument display}
\end{figure}

A child information is an information regarding a child of this issue point.
The first exclamation mark $!$ indicates that the following is an information regarding the rule which guided the issue point and relation with the TopRule of children for this issue point. 
Three kinds of relations exists:
$>$ indicates that the issue point has a stronger TopRule relation than all children.
$!$ indicates that it is either strong or unrelated.
$<$ indicates that it is weaker than at least one of the children.

The first figure of $1(0)$ shows the number of a child.
The figure inside the parentheses show the number of the child which a debate has ended.
The debate is considered to have ended when either of the following occurs;
the prosecutor and defendant both have noticed their argument to their
opponent by {\sl notice} command,
the opponent have noticed that the debate has finished by the {\sl finish} command. 

An intermediate proposition is a conclusion department of a rule.
A blank before an intermediate proposition shows that the one without an indentation was guided by the one with the indentation.

\subsubsection{Argument list display of a viewpoint}\label{Argument list display of a viewpoint}

This section explains the form to display multiple arguments displayed in a counter argument list display window and argument list display window.
The following is an example of a counter argument list display:

\begin{figure}[htb]
%\begin{center}
\begin{verbatim}
  Arg ID: p0
  Issue Pos: 13
  Consequence: with_criminal_intent(a_object=person
  RuleName: penal235otsuka_theory
  Num Of Counter: 1
  0 : *  !  >:0!:1<:0 -with_criminal_intent(a_object=person
\end{verbatim}
%\end{center}
\caption{Example of a counter argument list display}\label{Example of a counter argument list display}
\end{figure}

\begin{description}
\item[Arg identification:] Parent argument identification.
\item[Issue Pos:] Issue point number of a parent argument.\\
This child list display shows the issue point of parent argument.
\item[Consequence:] THEN part of the rule that guides an issue point corresponding to the issue point number of a parent argument.
\item[RuleName:] Name of the rule that guides an issue point corresponding to the issue point number of a parent argument.
\item[Num Of Counter:] Number of counter arguments to an issue point corresponding to the issue point number of a parent argument.
\item[0:] a child number.\\
\verb+*+ indicates that a counter argument has been completed by a {\sl pose} command.\\
\verb+!+ shows the TopRule relation between the issue point of the parent argument and argument corresponding to the issue point, and \verb+>+ means the parent is weaker than, \verb+!+ means unrelated to, and \verb+<+ means stronger than the child.\\
\verb+>:0!:1<:0+ shows the number of strong(\verb+>+), unrelated(\verb+!+), and weak(\verb+<+).
\end{description}

\subsubsection{To see an argument to counter argue}

Select Backward.
The argument of the opponent to counter argue is displayed in the show argument area.

\subsubsection{To see a counter argument}

Choose by double clicking an intermediate proposition as an issue point from the arguments currently displayed.
A list of counter arguments which regards the intermediate proposition as an issue point will be displayed.

\begin{figure}[htb]
\begin{center}
\epsfile{file=./cargwin.ps,height=6cm}
\end{center}
\caption{Counter argument selecting window}
\end{figure}

A list of a conclusion of counter arguments is displayed in this window.
Double click on a conclusion of the counter argument to see.
The counter argument selecting window closes and the selected argument is displayed in the show argument area.
The {\sl Forward} button may also be used to see a counter argument.
In this case the selecting operation of a counter argument is omitted.
The argument displayed before using the {\sl Backward} button is displayed in the show argument area.
Generally, the {\sl Forward} button is used to return to the original
argument after using the {\sl Backward} button.

\subsubsection{To see other counter arguments on the same issue point}

Select either the {\sl Next} button or {\sl Previous} button to display other arguments debating on the same issue point in the show argument area.
Since the display order of the {\sl Next} button and {\sl Previous} button are opposite, one can be used to return to the previous display.

\subsubsection{To see the argument that is a focus of a debate}

Select {\sl Focus}.

\subsubsection{To display the argument with the argument identification known}

Select {\sl Jump} to open the argument identification specifying window.
Enter the argument identification and hit return.
The argument of the entered argument identification will be displayed.  Notice that argument identification exists only for the argument which have been claimed.

\subsubsection{To retrieve a past claim}

Select {\sl ClaimList} to display the claim list window.

\begin{figure}[htb]
\begin{center}
\epsfile{file=./claimwin.ps,width=15cm}
\end{center}
\caption{Claim list window}\label{Claim list window}
\end{figure}

Displayed at the left is the claim of a prosecutor and the right is the claim of a defendant.
The display form is similar to the argument list display. (Refer to '\ref{Argument list display of a viewpoint} Argument list display of a viewpoint')
Double click the argument to display.

\subsubsection{To change the argument display into a natural language}

Select {\sl Form} to display the display form selecting window.

\begin{figure}[htb]
\begin{center}
\epsfile{file=./formwin.ps,height=3cm}
\end{center}
\caption{Display form selecting window}\label{Display form selecting window}
\end{figure}

To display the argument display in H term format, select {\sl H-term}.  To display the argument display in a natural language, select {\sl Natural Language}.  In this case, the Explain information added to the rule is used.

\section{ Operation in DRS window}

\begin{quotation}
The function of DRS(Defeasible Reasoning System) is to analyze problem in one viewpoint,
and generates all arguments and all its counter arguments.

DRS window made from initial window.

DRS window is same to remove operation of debate from debate window.

Therefore method of setting type,viewpoint,goal,and rule is refered
to 'Operation in the debate window'.

It is necessary for openning DRS window to finish load of type,viewpoint goal and rule.

Can't open DRS window and debate window in same time.

\end{quotation}

\begin{figure}[htb]
\begin{center}
\epsfile{file=./drswin.ps,height=12cm}
\end{center}
\caption{DRS window}\label{DRS window}
\end{figure}

\begin{itemize}

\item DRS setting part 
\begin{description}
\item[View button:] Sets a viewpoin.
\item[Unit button:] Sets a use knowledge.
\item[Goal button:] Sets a goal.
\item[Solve button:] Generates arguments.
\item[Quit  button:] Closes a DRS window.
\end{description}

\item DRS setting browser part
\begin{description}
\item[Goal area:] Displays the current goal.
\item[View area:] Displays the current viewpoint.
\end{description}

\item DRS state browser part
\begin{description}
\item[Show argument area:] Displays arguments.
Displays a Unified conclusion (intermediate proposition) of rules with an inentation, depending on the reference relation.
Each rule composes the proof tree of an argument.
\end{description}

\item{Argument browser operation part}
\begin{description}
\item[Backward button:] Displays the argument which the counter argument currently displayed in the show argument area is trying to override.
\item[Forward button:] Returns to the argument displayed immediately
  before selecting {\sl Backward}, in the case where {\sl Backward} was used.
Returns to one of the counter arguments to the argument displayed 
 in the show argument area.
\item[Next button:] Advances to the next counter argument, if it exists, disputing on the same issue point number as the argument currently displayed in the show argument area.
\item[Previous button:] Advances to the previous counter argument, if it exists, disputing on the same issue point number as the argument currently displayed in the show argument area.
Displays in the opposite order as the {\sl Next} button.
\item[Explain button:] Explains the argument currently displayed in a natural language.
\item[Form button:] Specifies the display form of the intermediate proposition displayed in the show argument area.
There are two display formats available; one is an expression by an H term variable and the other is an explanation in a natural language.
\end{description}

\end{itemize}

\subsection{To open a DRS window}
Select {\sl DRS} in the initial window.
However, loading of required files must be completed for
 the {\sl DRS} button to become sensitive.

\subsection{To close a DRS window}
Select {\sl Quit} button.

\subsection{To generates arguments}

Select {\sl Solve} in the DRS window.
The viewpoint and use knowledge and the goal must be set.
Result of inference is displayed in Show argument area.

Operation of argument browser is refered to 'Argument browser operation part'.


\section{Attention}

Current HELIC-II is not supported as follows.

Display argument by using natural language.

Function of browser to see claimed arguments after debate.


\end{document}
